\documentclass[../../main.tex]{subfiles}
\graphicspath{{\subfix{../images/}}} % 指定图片目录，后续可以直接使用图片文件名。
\begin{document}
\section{2025年10月博士生资格考试注解}

\subsection{流体力学}
Poiseuille 定律. 
\begin{enumerate}
    \item \textbf{列举涉及的重要物理量及其量纲}
    
    \begin{table}[h]
    \centering
    \begin{tabular}{|c|c|c|c|c|c|}
    \hline
    物理量 & 压强差 & 管长 & 半径 & 动力粘度 & 体积流量 \\
    \hline
    符号 & $\Delta P$ & $L$ & $R$ & $\mu$ & $Q$ \\
    \hline
    量纲 & $\mathrm{ML^{-1}T^{-2}}$ & $\mathrm{L}$ & $\mathrm{L}$ & $\mathrm{ML^{-1}T^{-1}}$ & $\mathrm{L^{3}T^{-1}}$ \\
    \hline
    \end{tabular}
    \end{table}
    
    [HINT] 剪切应力 $\tau = \mu\times$ 速度梯度($\begin{aligned}\frac{\mathrm{d}v}{\mathrm{d}r}\end{aligned}$) 
    
    \item \textbf{通过量纲分析流量 $Q\propto \Delta P$ [提示: $Q\propto \frac{\Delta P}{L}$]}
    
    设 $\begin{aligned}Q \sim \left(\frac{\Delta P}{L}\right)^{a} R^{b} \mu^{c}\end{aligned}$, 而根据第一问中, 

    \begin{table}[h]
\centering
\begin{tabular}{|c|c|c|c|c|}
\hline
 & 体积流量 & 压强梯度 & 半径 & 动力粘度 \\
\hline
符号 & $Q$ & $\frac{\Delta P}{L}$ & $R$ & $\mu$ \\
\hline
量纲 & $\mathrm{L^{3}T^{-1}}$ & $\mathrm{ML^{-2}T^{-2}}$ & $\mathrm{L}$ & $\mathrm{ML^{-1}T^{-1}}$ \\
\hline
\end{tabular}
\end{table}

这实际形成了一个待解的线性方程组

$$
\begin{aligned}
\mathrm{L^{3}T^{-1}} = \mathrm{(ML^{-2}T^{-2})}^{a} \cdot \mathrm{L}^{b} \cdot \mathrm{(ML^{-1}T^{-1})}^{c} = \mathrm{M}^{a+c} \mathrm{L}^{-2a+b-c} \mathrm{T}^{-2a-c}
\end{aligned}
$$

指数比较:

$$
\begin{cases}
a + c &= 0 \\
-2a + b - c &= 3 \\
-2a - c &= -1
\end{cases}
$$

解得 $a = 1, b = 4, c = -1$. 因此得到 Poiseuille 定律的量纲形式

$$
Q \propto \frac{\Delta P}{L}  \frac{R^{4}}{\mu}
$$

    \item \textbf{写出约化 navier-stokes 方程, 证明只需求解速度分布 $v(r)$}
    
一般的 Navier-Stokes 方程为

$$
\rho \left(\frac{\partial\vec{v}}{\partial t} + \vec{v}\cdot\nabla\vec{v}\right) = -\nabla P + \mu\nabla^{2}\vec{v}
$$

对于圆柱管中的层流, 可以做出以下简化: 1. 稳态: $\frac{\partial}{\partial t} = 0$; 2. 层流: 湍流项 $\vec{v}\cdot\nabla\vec{v} = 0$; 3. 轴对称: $\vec{v} = v(r)\hat{z}$; 4. 不可压缩流体: $\nabla\cdot\vec{v} = 0$.

因此得到简化后的 Navier-Stokes 方程

$$
0 = -\frac{\mathrm{d}P}{\mathrm{d}z} + \mu\nabla^{2}v
$$

柱坐标下, 

$$
\nabla^{2} = {\color{red}{\frac{1}{r}\frac{\mathrm{d}}{\mathrm{d}r}\left(r\frac{\mathrm{d}}{\mathrm{d}r}\right)}} + \frac{1}{r^{2}}\frac{\mathrm{d}^{2}}{\mathrm{d}\theta^{2}} + \frac{\mathrm{d}^{2}}{\mathrm{d}z^{2}}
$$

而该模型下 $\vec{v} = v({\color{red}{r}})\hat{z}$, 因此只需关注红色项.

    \item \textbf{代入边界条件精确解 $v(r)$}
    
边界条件为 1. 中心不发散: $C_{1} = 0$; 2. 管壁无滑移: $v(R) = 0$. 因此

$$
0 = -\frac{\Delta P}{4\mu L}R^{2} + C_{2} \Rightarrow C_{2} = \frac{\Delta P}{4\mu L}R^{2}
\Rightarrow v(r) = \frac{\Delta P}{4\mu L}(R^{2} - r^{2})
$$

    \item \textbf{精确解流量 $Q$}
    
$$
\begin{aligned}
Q &= \int v\mathrm{d}A = \int_{0}^{R}v(r) 2\pi r\mathrm{d}r = \int_{0}^{R} \frac{\Delta P}{4\mu L}(R^{2} - r^{2}) 2\pi r\mathrm{d}r\\
&= \frac{\Delta P}{4\mu L}2\pi \int_{0}^{R} (R^{2}r - r^{3})\mathrm{d}r = \frac{\Delta P}{4\mu L}2\pi (R^{2}\frac{r^{2}}{2} - \frac{1}{4}r^{4})\Big|_{0}^{R}\\
&= \frac{\Delta P}{4\mu L}2\pi \frac{1}{4}R^{4} = \frac{\pi R^{4}}{8\mu }\frac{\Delta P}{L}
\end{aligned}
$$

    \item \textbf{比较精确解和量纲分析结果的异同}

    \begin{itemize}
        \item 量纲分析结果: $\begin{aligned}Q \propto \frac{\Delta P}{L}  \frac{R^{4}}{\mu}\end{aligned}$;
        \item 精确解: $\begin{aligned}Q = \frac{\pi R^{4}}{8\mu }\frac{\Delta P}{L}\end{aligned}$
    \end{itemize}

量纲分析得到的解形式和精确解的完全一致, 而精确解则是得出了具体的系数 $\pi/8$. 
\end{enumerate}

\subsection{应力学}

\begin{enumerate}
    \item \textbf{画图说明应力概念}
    
应力被定义为 "单位面积上的内力", 因此量纲与压强相同: $\begin{aligned}\vec{t} = \lim_{\Delta A\to 0}\frac{\Delta\vec{F}}{\Delta A}\end{aligned}$. 显然 $\vec{t}$ 是和面元 $\Delta A$ 的选取有关的, 后者通常通过其法向量 $\hat{n}$ 和所在位置 $(x,y,z)$ 定义. 

接下来研究 $\Delta\vec{F}$ 的方向. 不妨设 $\hat{n}$ 为 $x$ 轴正方向, 则

$$
\begin{aligned}
\Delta\vec{F} &= \Delta F_{x}\hat{x} + (\Delta F_{y}\hat{y} + \Delta F_{z}\hat{z})\\
&= \Delta F_{x}\hat{n} + (\Delta F_{y}\hat{y} + \Delta F_{z}\hat{z})
\end{aligned}
$$

其中 $\Delta F_{x}\hat{n}$ 为法向分量(normal stress), $\Delta F_{y}\hat{y} + \Delta F_{z}\hat{z}$ 为切向分量(shear stress). 

可以将 $\vec{t}_{x}$ (下标 $x$ 表示将 $\hat{n}$ 选取为 $\hat{x}$) 写作这样的形式: 

$$
\begin{aligned}
\vec{t}_{x} &= \lim_{\Delta A\to 0}\frac{\Delta F_{x}}{\Delta A}\hat{x} + \lim_{\Delta A\to 0}\frac{\Delta F_{y}}{\Delta A}\hat{y} + \lim_{\Delta A\to 0}\frac{\Delta F_{z}}{\Delta A}\hat{z}\\
&= \sigma_{x}\hat{x} + \tau_{xy}\hat{y} + \tau_{xz}\hat{z}
\end{aligned}
$$

其中 $\sigma_{x}$ 也可写作 $\sigma_{xx}$ 的形式. $\sigma_{ij}$ 即可理解为, 施加在法线为 $i$ 轴的面上, 沿着 $j$ 轴的应力分量大小. 


按照这个规则, 写出 $\hat{n}$ 为 $\hat{y}$ 和 $\hat{z}$ 时的应力: 

$$
\begin{aligned}
\vec{t}_{y} &= \tau_{yx}\hat{x} + \sigma_{y}\hat{y} + \tau_{yz}\hat{z}\\
\vec{t}_{z} &= \tau_{zx}\hat{x} + \tau_{zy}\hat{y} + \sigma_{z}\hat{z}
\end{aligned}
$$

若将所有应力分量统一记为 $\sigma_{ij}$, 则可得到矩阵:

$$
\mathbf{\sigma} = \begin{bmatrix}
\sigma_{xx} & \sigma_{xy} & \sigma_{xz}\\
\sigma_{yx} & \sigma_{yy} & \sigma_{yz}\\
\sigma_{zx} & \sigma_{zy} & \sigma_{zz}
\end{bmatrix}
$$

因为物体内部需要满足力矩平衡, 因此这是一个对称矩阵, 即满足

$$
\sigma_{ij} = \sigma_{ji}
$$

    \item \textbf{解释为什么应力是张量}
    
应力和面元的选取有关: 

$$
\vec{t}_{i} = \sum_{j}\sigma_{ij}\hat{n}_{j}
$$

要求应力同时蕴含面元选取(面法线 $i$ 轴) 和应力分量( $j$ 轴)的信息, 因此应力必然会形成一个 $3\times 3$ 的二阶张量(矩阵). 


    \item \textbf{说明什么时候应力可被简化为压强}

    压强被定义为

$$
p = \frac{F_{n}}{A}
$$

因此只沿面元法向. 既然面元的选取是任意的, 因此压强具有各向同性. 对于一般的流体, 其应力张量为

$$
\sigma_{ij} = -p\delta_{ij} + \tau_{ij}
$$

其中 $\tau_{ij}$ 表示剪切应力, 通常是液体的粘性导致的. 因此对于静止流体具有

$$
\mathbf{\sigma} = -p\mathbb{I} = \begin{bmatrix}
-p & 0 & 0\\
0 & -p & 0\\
0 & 0 & -p
\end{bmatrix}
$$

因此在各向同性/无剪切的情况下, 应力可以被简化为压强. 
\end{enumerate}

\subsection{材料学}

\begin{enumerate}
    \item \textbf{解释概念 "杨氏模量"}
\begin{enumerate}
    \item 应力(stress): $\begin{aligned}\sigma = \frac{F}{A}\end{aligned}$, 即单位面积上承受的力, 量纲和压强相同;
    \item 应变(strain): $\begin{aligned}\epsilon = \frac{\Delta L}{L}\end{aligned}$, 即长度的相对变化
\end{enumerate}

杨氏模量(Young's modulus)定义为应力与应变的比值:

$$
E = \frac{\sigma}{\epsilon}
$$

    \item \textbf{排序花岗岩, 钢铁, 橡胶材料的杨氏模量大小}

$E$ 越大/小, 材料越硬/软.

\begin{table}[h]
\centering
\begin{tabular}{|c|c|}
\hline
材料 & 杨氏模量 $E$ \\
\hline
橡胶 & $10^{6}\sim 10^{7}\,\mathrm{Pa}$ \\
花岗岩 & $30\sim 70\,\mathrm{GPa}$ \\
钢 & $200\,\mathrm{GPa}$ \\
\hline
\end{tabular}
\end{table}

$\boxed{E_{\text{橡胶}} < E_{\text{花岗岩}} < E_{\text{钢}}}$.

橡胶是高分子链, 具有构型熵弹性; 花岗岩是离子键和共价键混合, 刚度高但较脆; 钢是金属键, 刚度高且延展性更好. 

    \item \textbf{解释概念 "泊松比", 并且举例说明负泊松比材料}

\begin{itemize}
    \item 轴向应变: $\epsilon_{\parallel}$, 在受力方向上的相对长度变化;
    \item 横向应变: $\epsilon_{\perp}$, 垂直于受力方向上的相对长度变化.
\end{itemize}

泊松比(Poisson's ratio)定义为横向应变($\epsilon_{\perp}$)与轴向应变($\epsilon_{\parallel}$)的负比值: $\begin{aligned}
    \nu = -\frac{\epsilon_{\perp}}{\epsilon_{\parallel}}
\end{aligned}$. 

一般的材料都是如同拉面般, 受拉伸时变细而受压缩时变粗, 即 $\nu > 0$. 

单轴拉伸下, 泊松比和体积变化的关系: $\begin{aligned}
    \frac{\Delta V}{V} = (1 - 2\nu)\epsilon
\end{aligned}$

负泊松比($\nu<0$) 则是指的受拉伸反而变粗, 受压缩反而变细的材料. 这类材料是因为内部具有独特的微观结构, 使得内部的结构单元在受力时产生旋转或者变形, 如 re-entrant honeycomb 结构.

\begin{figure}[htbp]
\centering
\includegraphics[width=0.4\textwidth]{image/-nu.png}
\end{figure}

\end{enumerate}

\subsection{布朗运动}

\begin{enumerate}
    \item \textbf{解释概念 "布朗运动"}

布朗运动(Brownian motion)指的是悬浮在液体或气体中的微小颗粒, 由于受到周围分子热运动的不断随机撞击而产生的无规则运动. 

均方位移:

$$
\langle x^{2}(t)\rangle = 2Dt, \quad D = \frac{k_{B}T}{\gamma}
$$
    \item \textbf{写出朗之万方程, 并对其进行变量解释}
    
$$
m\frac{\mathrm{d}v}{\mathrm{d}t} = -\gamma v + \xi(t)
$$

$m$: 粒子质量; $-\gamma v$: 粘滞阻力; $\xi(t)$: 随机力, 且满足

$$
\begin{aligned}
\langle \xi(t)\rangle &= 0\\ \langle \xi(t)\xi(t^{\prime})\rangle &= 2\gamma k_{B}T \delta(t-t^{\prime})
\end{aligned}
$$
    \item \textbf{解释概念 "涨落耗散关系"}

$$
\boxed{\langle \xi(t)\xi(t^{\prime})\rangle = 2\gamma k_{B}T \delta(t-t^{\prime})}
$$

    \item \textbf{假定布朗粒子处于重力势中, 求解其在高度上的密度分布}

设竖直向上为正方向, 则朗之万方程形式为

$$
m\frac{\mathrm{d}^{2}h}{\mathrm{d}t^{2}} = -\gamma \frac{\mathrm{d}h}{\mathrm{d}t} - mg + \xi(t)
$$

在该模型中, 有意义的解一般考虑布朗粒子处于过阻尼极限(粘滞远大于惯性), 则方程改写为

$$
0 = -\gamma\frac{\mathrm{d}h}{\mathrm{d}t} - mg + \xi(t)\\
\Rightarrow \frac{\mathrm{d}h}{\mathrm{d}t} = -\frac{mg}{\gamma} + \frac{\xi(t)}{\gamma}
$$

若有微分方程

$$
\frac{\mathrm{d}h}{\mathrm{d}t} = A(h) + \eta(t)
$$

其中

$$
\langle \eta(t)\eta(t^{\prime})\rangle = 2D\delta(t-t^{\prime})
$$

则概率密度可改写为 Fokker-Planck 方程

$$
\frac{\partial f}{\partial t} = - \frac{\partial (Af)}{\partial h} + D\frac{\partial^{2}f}{\partial h^{2}}
$$

在该模型中, 

$$
A = -\frac{mg}{\gamma},\quad D = \frac{k_{B}T}{\gamma}
$$

因此 Fokker-Planck 方程为

$$
\frac{\partial f}{\partial t} = \frac{mg}{\gamma} \frac{\partial f}{\partial h} + \frac{k_{B}T}{\gamma}\frac{\partial^{2}f}{\partial h^{2}}
$$

在稳态/热平衡下, $\begin{aligned}\frac{\partial f}{\partial t} = 0\end{aligned}$, 即得到 ODE

$$
0 = \frac{mg}{\gamma} \frac{\mathrm{d}f}{\mathrm{d}h} + \frac{k_{B}T}{\gamma}\frac{\mathrm{d}^{2}f}{\mathrm{d}h^{2}}
$$

设辅助函数 $\begin{aligned}u = \frac{\mathrm{d}f}{\mathrm{d}h}\end{aligned}$, 则方程为

$$
\frac{\mathrm{d}u}{\mathrm{d}h} = -\frac{mg}{k_{B}T}u
\Rightarrow \frac{\mathrm{d}u}{u} = -\frac{mg}{k_{B}T}\mathrm{d}h
\Rightarrow u = Ce^{-\frac{mg}{k_{B}T}h}
$$

因此

$$
f(h) = -\frac{Ck_{B}T}{mg}\exp{\left(-\frac{mg}{k_{B}T}h\right)} + D
$$

代入边界条件 $\begin{aligned}\lim_{h\to\infty}f = 0\end{aligned}$, 得到 $D = 0$. 

归一化条件 $\begin{aligned}\int_{0}^{\infty}f(h)\mathrm{d}h = 1\end{aligned}$ 要求解形式为 $\begin{aligned}f(h) = \frac{1}{h_{0}}e^{-h/h_{0}}\end{aligned}$, 其中 $\begin{aligned}h_{0} = \frac{k_{B}T}{mg}\end{aligned}$.

因此

$$
\boxed{f(h) = \frac{mg}{k_{B}T}e^{-\frac{mg}{k_{B}T}h}}
$$

\begin{center}
\rule{0.8\textwidth}{0.4pt}
\end{center}

{\color{gray}{[Discussion] Fokker-Planck 方程的导出. 

一维 Ito SDE(stochastic differential equation, 随机微分方程)为 

$$
\mathrm{d}x = A(x)\mathrm{d}t + \sqrt{2D}\mathrm{d}W_{t}
$$

> $\langle\mathrm{d}W_{t}\rangle = 0$, $\langle(\mathrm{d}W_{t})^{2}\rangle = \mathrm{d}t$. 

设概率密度为 $f(x,t)$, 且引入辅助函数 $\phi(x)$, 其期望为

$$
\langle \phi(x)\rangle = \int \phi(x)f(x,t)\mathrm{d}x
$$

微元展开:

$$
\mathrm{d}\phi = \phi^{\prime}(x)\mathrm{d}x + \frac{\phi^{\prime\prime}(x)}{2}(\mathrm{d}x)^{2}
$$

代入 $\mathrm{d}x = A(x)\mathrm{d}t + \sqrt{2D}\mathrm{d}W$, 需要注意到性质

$$
W_{t+\Delta t} - W_{t}\sim \mathcal{N}(0,\Delta t)\Rightarrow \Delta W\sim\sqrt{\Delta t}
$$

因此有 

$$
(\mathrm{d}x)^{2} = (A(x)\mathrm{d}t + \sqrt{2D}\mathrm{d}W)^{2} = 2D\mathrm{d}t
$$

于是

$$
\begin{aligned}
\mathrm{d}\phi &= \phi^{\prime}(x)A\mathrm{d}t + \phi^{\prime}\sqrt{2D}\mathrm{d}W + D\phi^{\prime\prime}\mathrm{d}t\\
&= \left[A\phi^{\prime} + D\phi^{\prime\prime}\right]\mathrm{d}t + \sqrt{2D}\phi^{\prime}\mathrm{d}W
\end{aligned}
$$

对两侧取期望, 且注意到 $\langle \mathrm{d}W\rangle = 0$, 因此

$$
\mathrm{d}\langle\phi\rangle = \langle A\phi^{\prime} + D\phi^{\prime\prime}\rangle \mathrm{d}t \Rightarrow \frac{\mathrm{d}\langle\phi\rangle}{\mathrm{d}t} = \langle A\phi^{\prime} + D\phi^{\prime\prime}\rangle
$$

将 $\langle \cdot\rangle$ 展开为积分形式 $\int \cdot f\mathrm{d}x$: 

$$
\frac{\mathrm{d}}{\mathrm{d}t} \int\phi f\mathrm{d}x = \int (A\phi^{\prime} + D\phi^{\prime\prime})f\mathrm{d}x
$$

使用分部积分: 

$$
\begin{aligned}
\int A\phi^{\prime}f\mathrm{d}x &= \int Af \mathrm{d}\phi = Af\phi - \int \phi \frac{\partial(Af)}{\partial x}\mathrm{d}x\\
\int D\phi^{\prime\prime}f\mathrm{d}x &= \int Df\mathrm{d}\phi^{\prime} = Df\phi^{\prime} - \int \phi^{\prime}\frac{\partial Df}{\partial x}\mathrm{d}x\\
&= Df\phi^{\prime} - \int \frac{\partial Df}{\partial x}\mathrm{d}\phi = Df\phi^{\prime} - \frac{\partial D f}{\partial x}\phi + \int \phi \frac{\partial^{2}Df}{\partial x^{2}}\mathrm{d}x
\end{aligned}
$$

假定边界条件为 $0$, 因此化为形式

$$
\int A\phi^{\prime}f\mathrm{d}x = -\int \phi \frac{\partial }{\partial x}(Af)\mathrm{d}x\\
\int D\phi^{\prime\prime}f\mathrm{d}x = \int \phi D\frac{\partial^{2}f}{\partial x^{2}}\mathrm{d}x
$$

代回原等式中:

$$
\frac{\mathrm{d}}{\mathrm{d}t}\int\phi f\mathrm{d}x = \int \phi \left[
    -\frac{\partial }{\partial x}(Af) + D\frac{\partial^{2}f}{\partial x^{2}}
\right]\mathrm{d}x
$$

因为对任意 $\phi$ 都成立, 因此

$$
\frac{\mathrm{d}}{\mathrm{d}t}\int\phi f\mathrm{d}x = \int \phi \frac{\partial f}{\partial t}\mathrm{d}x
$$

得到 Fokker-Planck 方程

$$
\boxed{\frac{\partial f}{\partial t} = -\frac{\partial }{\partial x}(Af) + D\frac{\partial^{2}f}{\partial x^{2}}}
$$}}

\end{enumerate}

\subsection{聚合物}

\begin{enumerate}
    \item \textbf{解释概念 "理想链" 和 "非理想链", 并阐述两者之间的区别}
    
\begin{itemize}
    \item 理想链(Ideal chain)

忽略单元间的相互作用而可近似为 random walk 的模型. 设 $N$ 个等长链段 $b$: 

$$
\begin{aligned}
\vec{R} &= \sum_{i=1}^{N}\vec{b}_{i},\quad \langle\vec{b}_{i}\cdot\vec{b}_{j}\rangle = b^{2}\delta_{ij}\\
\vec{R}^{2} &= \vec{R}\cdot\vec{R} = \left(\sum_{i=1}^{N}\vec{b}_{i}\right)\cdot\left(\sum_{j=1}^{N}\vec{b}_{j}\right) = \sum_{i=1}^{N}\sum_{j=1}^{N}\vec{b}_{i}\cdot\vec{b}_{j}\\
\langle R^{2}\rangle &= \sum_{i=1}^{N}\sum_{j=1}^{N}\langle \vec{b}_{i}\cdot\vec{b}_{j}\rangle = \sum_{i=1}^{N}\sum_{j=1}^{N} b^{2}\delta_{ij} = Nb^{2}
\end{aligned}
$$

因此 $R\sim N^{1/2}$.
    \item 非理想链(Non-ideal chain)

单元存在排斥体积(self-avoidance), 单元之间不能重叠, 存在相互排斥作用. 因此相较于理想链, 半径将会修正为

$$
R\sim bN^{\nu}
$$

其中 $\nu$ 被称作 Flory 指数, 其具体大小和 Polymer 所处的溶剂环境有关. 
\end{itemize}

    \item \textbf{求解三维空间中理想链的熵弹性力}
    
{\color{gray}{[Recall] 一维空间中的 Random walk

首先考虑最简单的 1D 模型. 设步长为 $a$, 总步数为 $N = n_{l} + n_{r}$, 则最后位置为 $R = (n_{r}-n_{l})a$ (假定向右为正方向). 向左和向右的概率均为 $1/2$, 因此处于位置 $R$ 的概率为

$$
\begin{aligned}
p(R,N) &= \frac{N!}{n_{l}!n_{r}!}\frac{1}{2^{N}} = \frac{N!}{(\frac{Na+R}{2})!(\frac{Na-R}{2})!}\frac{1}{2^{N}}\\
\ln{p} &\approx \ln{\sqrt{\frac{2}{\pi N}}} - \frac{R^{2}}{2Na^{2}}\\
p(R,N) &\approx \sqrt{\frac{2}{\pi N}}\exp{\left(-\frac{R^{2}}{2Na^{2}}\right)}
\end{aligned}
$$

计算 $\begin{aligned}\int_{-Na}^{Na}p\mathrm{d}R = 2a\end{aligned}$, 因此概率密度需要再归一化:

$$
p(R,N) = \frac{1}{\sqrt{2\pi Na^{2}}}\exp{\left(-\frac{R^{2}}{2Na^{2}}\right)}
$$

这是一个高斯分布, 均值为 $\bar{R} = 0$, 方差为 $\sigma_{R}^{2} = Na^{2}$. 

研究 Ideal chain 的熵弹性. 熵变被定义为 Boltzmann 熵差

$$
\Delta S = k_{B}\ln{\Omega(N,x)} - k_{B}\ln{\Omega(N,0)} = k_{B}\ln{\frac{p(N,x)}{p(N,0)}} = -\frac{1}{2}\frac{k_{B}}{Na^{2}}x^{2}
$$

自由能变 $\begin{aligned}
    \Delta G = \Delta U - T\Delta S = \frac{1}{2}\frac{k_{B}T}{Na^{2}}x^{2}
\end{aligned}$, 因此熵弹性力为 $\begin{aligned}
    f = \frac{\partial G(N,x)}{\partial x} = \frac{k_{B}T}{Na^{2}}x
\end{aligned}$. }}

\begin{center}
\rule{0.8\textwidth}{0.4pt}
\end{center}

推广至 3D 模型时, 则可将其视为三个独立的 1D 模型的叠加:

$$
\vec{R} = (R_{x}, R_{y}, R_{z}),\quad p(\vec{R}) = p(R_{x})p(R_{y})p(R_{z})
$$

而已得 1D 形式为 $\begin{aligned}
    p(R_{i},N_{i}) = \frac{1}{\sqrt{2\pi N_{i}a^{2}}}\exp{\left(-\frac{R_{i}^{2}}{2N_{i}a^{2}}\right)}
\end{aligned}$. 

由于 $N$ 充分大, 我们可以近似 $N_{x} = N_{y} = N_{z} = N/3$, 因此

$$
p(\vec{R},N) = p\left(R_{x}, \frac{N}{3}\right)p\left(R_{y}, \frac{N}{3}\right)p\left(R_{z}, \frac{N}{3}\right) = \left(\frac{3}{2\pi Na^{2}}\right)^{3/2}\exp{\left(-\frac{3R^{2}}{2Na^{2}}\right)}
$$

类似地, 熵变为 Boltzmann 熵差 $\begin{aligned}
    \Delta S = k_{B}\ln{\frac{p\left(\vec{R},N\right)}{p\left(\vec{0},N\right)}} = -\frac{3}{2}\frac{k_{B}}{Na^{2}}R^{2}
\end{aligned}$, 自由能为 $F = U - TS$, 因为 $(U, T, P)$ 不变, 所以能变 $\begin{aligned}
    \Delta F = - T\Delta S = \frac{3}{2}\frac{k_{B}T}{Na^{2}}R^{2}
\end{aligned}$, 熵弹性力 $\begin{aligned}
    f = 3 \frac{k_{B}T}{Na^{2}}R
\end{aligned}$. 

    \item \textbf{解释概念 "Flory 理论"}

描述 non-ideal chain 的理论. 相比 ideal chain 引入了单元间的体积作用, 因此自由能中存在 "熵弹性" 和 "体积排斥" 两部分的竞争. 

\begin{itemize}
    \item 熵弹性(entropic elasticity): $\begin{aligned}
        F_{\text{elastic}} = \frac{3}{2}\frac{k_{B}T}{Na^{2}}R^{2}
    \end{aligned}$

    \item 体积排斥(excluded volume): 单元的数密度为 $\begin{aligned}
        n\sim \frac{N}{R^{3}}
    \end{aligned}$, 引入排斥体积参数 $v$, 则体积排斥能密度 $\sim vn^{2}$. 将能密度在空间上积分, 则总共的体积排斥能写作 $\begin{aligned}
        F_{\text{int}}\sim k_{B}T vn^{2}R^{3} \sim k_{B}T v\frac{N^{2}}{R^{3}}
    \end{aligned}$
\end{itemize}

因此总的自由能为两者之和: $\begin{aligned}
    F(R) \sim k_{B}T \left(
    \frac{R^{2}}{Na^{2}} + v\frac{N^{2}}{R^{3}}
\right)
\end{aligned}$. 

为了使得自由能最小, 熵弹性倾向于令 $R\downarrow$, 而体积排斥倾向于令 $R\uparrow$. 求偏导以求得极小值: 

$$
\frac{\partial F}{\partial R} = k_{B}T \left(
    \frac{2R}{Na^{2}} - 3v\frac{N^{2}}{R^{4}}
\right) = 0 \Rightarrow R \sim aN^{3/5}
$$

可将上式进一步抽象为 $R\sim aN^{\nu}$, 其中 $\nu$ 被称作 Flory 指数. 需要注意的是不同维度下 Flory 指数也不同. 例如 1D 下 $\nu = 1$, 2D 下 $\nu = 3/4$. 3D 下 $\nu$ 的精确解约为 $\nu\approx 0.588$. 

    \item \textbf{若将聚合物溶解于某溶剂中, 阐述溶剂和温度对其平均半径的影响}

溶剂和温度会影响排斥体积参数 $\nu$. 根据单体和溶剂的亲和性, 可将溶剂进行分类: 

\begin{itemize}
    \item Good solvent(好溶剂): 单体和溶剂亲和, 从而体现出链单元间的排斥作用, $\nu>0$. 此时尺度为上面的 self-avoiding walk 解: $R\sim aN^{3/5}$;
    \item Theta solvent($\theta$ 溶剂): 吸引和排斥相互抵消, $\nu=0$. 在 $\theta$ 溶剂中聚合物表现为理想链(random walk): $R\sim aN^{1/2}$;
    \item Poor solvent(差溶剂): 单体和溶剂不亲和, 因此单体更倾向于相互接触, $\nu < 0$. 此时链会发生 "collapse(塌缩)" 现象, 即可将聚合物视为致密球体. 此时回过头检查自由能式, $\begin{aligned}
        F(R)\sim k_{B}T \left(
    \frac{R^{2}}{Na^{2}} - |v|\frac{N^{2}}{R^{3}}
\right)
    \end{aligned}$, 

当 $R\downarrow$ 时, $F\rightarrow -\infty$, 这显然是荒谬的. 引入三体排斥项 $\begin{aligned}
    \iiint wn^{3}\mathrm{d}R
\end{aligned}$, 则 Flory 自由能为

$$
\begin{aligned}
F(R) \sim k_{B}T \left(
    \frac{R^{2}}{Na^{2}} - |v|\frac{N^{2}}{R^{3}} + w\frac{N^{3}}{R^{6}}
\right)\approx k_{B}T \left(
    - |v|\frac{N^{2}}{R^{3}} + w\frac{N^{3}}{R^{6}}
\right)
\end{aligned}
$$

因为聚合物塌缩为球状($R\rightarrow 0$), 因此忽略熵弹性项 $R^{2}$. 求偏导以求极小值: 

$$
\frac{\partial F}{\partial R} = k_{B}T \left(
    3|v|\frac{N^{2}}{R^{4}} - 6w\frac{N^{3}}{R^{7}}
\right) = 0\\
\Rightarrow R^{3}\sim \frac{2w}{|v|}N \Rightarrow R\sim N^{1/3}
$$

也可以理解为, 因为聚合物形成致密球体, 此时密度约为常数, 即 $\begin{aligned}
    \rho\sim\frac{N}{R^{3}}\approx \text{const}
\end{aligned}$, 从而推测 $R^{3}\propto N\Rightarrow R\propto N^{1/3}$.
\end{itemize}

{\color{gray}{\begin{center}
\rule{0.8\textwidth}{0.4pt}
\end{center}

[Discussion] 聚合物混合熵的推导. 

设两种聚合物混合, 分子链数分别为 $n_{1}$ 和 $n_{2}$, 聚合度分别为 $N_{1}$ 和 $N_{2}$. 混合物总格点数 $M = n_{1}N_{1} + n_{2}N_{2}$. 设该混合晶格的配位数为 $z$. 现在将聚合度为 $N$ 的链放入已经有 $x$ 个格点被占据的晶格中. 

\begin{itemize}
    \item 第 1 单体: 可以放在任意一个空位上, 因此选择数为 $M-x$;
    \item 第 2 单体: 只能放在第 1 单体的配位, 此时空间中空格点的平均概率为 $\begin{aligned}p = \frac{M-x}{M}\end{aligned}$. 因此选择数为 $\begin{aligned}zp = z\frac{M-x}{M}\end{aligned}$;
    \item 第 $3-N$ 单体: 由于自体积排斥, 后续可选的格点数为 $\begin{aligned}(z-1)p = (z-1)\frac{M-x}{M}\end{aligned}$;
\end{itemize}

因此, 放入该链的选择数 $\nu$ 为

$$
\nu = (M-x) \cdot \left[z\frac{M-x}{M}\right] \cdot \left[(z-1)\frac{M-x}{M}\right]^{N-2}\overset{z\approx z-1}{\simeq} \left[\frac{z-1}{M}\right]^{N-1}(M-x)^{N}
$$

证明近似 $\begin{aligned}(M-x)^{N}\approx\frac{(M-x)!}{(M-x-N)!}\end{aligned}$. 展开等式右边: 

$$
\begin{aligned}
\frac{(M-x)!}{(M-x-N)!} = (M-x)(M-x-1)(M-x-2)\cdots [M-x-(N-1)]\\
= (M-x)^{N}\left[
    1\cdot \left(1 - \frac{1}{M-x}\right)\left(1 - \frac{2}{M-x}\right) \cdots\left(1 - \frac{N-1}{M-x}\right)
\right]
\end{aligned}
$$

由于 $M\gg N$, 因此 $\begin{aligned}1 - \frac{k}{M-x}\to 1\end{aligned}$. 

若将 $M-x$ 视为整体 $M$, 该近似公式也可写作 $\begin{aligned}\frac{M!}{(M-N)!} \approx M^{N}\end{aligned}$.

采用近似 $\begin{aligned}(M-x)^{N}\approx\frac{(M-x)!}{(M-x-N)!}\end{aligned}$, 第 $i$ 条放入聚合度为 $N$ 链的方法数为

$$
\nu_{i} = \left[\frac{z-1}{M}\right]^{N-1}\frac{(M-x_{i})!}{(M-x_{i}-N)!}
$$

\begin{itemize}
    \item 分离状态 ( I-链和 II-链各自占据自己纯物质的晶格)

对于 I-链, 总晶格数为 $M_{1} = n_{1}M_{1}$, 需往空晶格(初始 $x_{1}=0$) 放入 $n_{1}$ 条 I-链. 被占格点数 $x_{i} = (i-1)N $. 引入链的不可区分性因子 $\begin{aligned}\frac{1}{n_{1}!}\end{aligned}$, 得到总方法数为

$$
\begin{aligned}
\Omega_{1} &= \frac{1}{n_{1}!}\prod_{i=1}^{n_{1}}\nu_{1,i} = \frac{1}{n_{1}!}\prod_{i=1}^{n_{1}}\left[\frac{z-1}{M_{1}}\right]^{N_{1}-1}\frac{[M-(i-1)N]!}{(M_{1}-iN_{1})!}\\
&= \frac{1}{n_{1}!}\left[\frac{z-1}{M_{1}}\right]^{n_{1}(N_{1}-1)}\prod_{i=1}^{n_{1}}\frac{[M-(i-1)N]!}{(M_{1}-iN_{1})!}\\
&= \frac{1}{n_{1}!}\left[\frac{z-1}{M_{1}}\right]^{n_{1}(N_{1}-1)} 
\frac{M_{1}!}{(M_{1}-N_{1})!}
\frac{(M_{1}-N_{1})!}{(M_{1}-2N_{1})!}\cdots
\frac{N_{1}!}{0!}\\
&= \frac{M_{1}!}{n_{1}!}\left[\frac{z-1}{M_{1}}\right]^{n_{1}(N_{1}-1)}
\end{aligned}
$$

II-链同理可得 $\begin{aligned}\Omega_{2} = \frac{M_{2}!}{n_{2}!}\left[\frac{z-1}{M_{2}}\right]^{n_{2}(N_{2}-1)}\end{aligned}$, 分离状态下的总状态数为 

$$
\begin{aligned}
\Omega_{\text{sep}} &= \Omega_{1}\Omega_{2} = \frac{M_{1}!}{n_{1}!}\left[\frac{z-1}{M_{1}}\right]^{n_{1}(N_{1}-1)}  \frac{M_{2}!}{n_{2}!}\left[\frac{z-1}{M_{2}}\right]^{n_{2}(N_{2}-1)}\\
&= \frac{M_{1}!M_{2}!}{n_{1}!n_{2}!} \left[\frac{z-1}{M_{1}}\right]^{n_{1}(N_{1}-1)} \left[\frac{z-1}{M_{2}}\right]^{n_{2}(N_{2}-1)}
\end{aligned}
$$
    \item 混合状态. 现在晶格总数为 $M = n_{1}N_{1} + n_{2}N_{2}$. 前面已经推导出, $M$ 总格点数中已有 $x$ 格点被占, 放入 $N$ 聚合度链的方法数 $\nu$ 为

$$
\nu(x,N) = \left[\frac{z-1}{M}\right]^{N-1}\frac{(M-x)!}{(M-x-N)!}
$$

\begin{enumerate}
    \item 填入 $n_{1}$ 条 I-链. 对于第 i 条被填入的 I-链, 已被占格点数 $x_{1,i}=(i-1)N_{1}$. 那么填入该 (I,i) 链的方法数为

$$
\nu_{1,i} = \left[\frac{z-1}{M}\right]^{N_{1}-1}\frac{[M-(i-1)N_{1}]!}{(M-iN_{1})!}
$$

引入 I-链不可区分性的因子 $\begin{aligned}\frac{1}{n_{1}!}\end{aligned}$, 填入这 $n_{1}$ 条 I-链的总方法数为

$$
\begin{aligned}
W_{1} &= \frac{1}{n_{1}!}\prod_{i=1}^{n_{1}}\nu_{1,i} = \frac{1}{n_{1}!}\prod_{i=1}^{n_{1}}  \left[\frac{z-1}{M}\right]^{N_{1}-1} \frac{[M-(i-1)N_{1}]!}{(M-iN_{1})!}\\
&= \frac{1}{n_{1}!} \left[\frac{z-1}{M}\right]^{n_{1}(N_{1}-1)} \prod_{i=1}^{n_{1}} \frac{[M-(i-1)N_{1}]!}{(M-iN_{1})!}\\
&= \frac{1}{n_{1}!} \left[\frac{z-1}{M}\right]^{n_{1}(N_{1}-1)} \frac{M!}{(M-N_{1})!} \frac{(M-N_{1})!}{(M-2N_{1})!}\cdots \frac{[M-(n_{1}-1)N_{1}]!}{(M-n_{1}N_{1})!}\\
&= \frac{1}{n_{1}!}\left[\frac{z-1}{M}\right]^{n_{1}(N_{1}-1)}\frac{M!}{(M-n_{1}N_{1})!}
\end{aligned}
$$
    \item 继续填入 $n_{2}$ 条 II-链. 此时 I-链已经占据 $n_{1}N_{1}$ 个格点. 所以对于放入的第 $j$ 条 II-链, 已被占据的格点数为 $x_{2,j} = n_{1}N_{1} + (j-1)N_{2}$, 填入 (II,j) 链的方法数为

$$
\nu_{2,j} = \left[\frac{z-1}{M}\right]^{N_{2}-1}\frac{[M-n_{1}N_{1}-(j-1)N_{2}]!}{(M-n_{1}N_{1}-jN_{2})!}
$$

引入 II-链的不可区分性因子 $\begin{aligned}\frac{1}{n_{2}!}\end{aligned}$, 那么填入这 $n_{2}$ 条 II-链的总方法数为

$$
\begin{aligned}
W_{2} &= \frac{1}{n_{2}!}\prod_{j=1}^{n_{2}}\nu_{2,j} = \frac{1}{n_{2}!}\prod_{j=1}^{n_{2}} \left[\frac{z-1}{M}\right]^{N_{2}-1}\frac{[M-n_{1}N_{1}-(j-1)N_{2}]!}{(M-n_{1}N_{1}-jN_{2})!}\\
&= \frac{1}{n_{2}!} \left[\frac{z-1}{M}\right]^{n_{2}(N_{2}-1)}\prod_{j=1}^{n_{2}}\frac{[M-n_{1}N_{1}-(j-1)N_{2}]!}{(M-n_{1}N_{1}-jN_{2})!}\\
&= \frac{1}{n_{2}!} \left[\frac{z-1}{M}\right]^{n_{2}(N_{2}-1)} \\
&\times \frac{(M-n_{1}N_{1})!}{(M-n_{1}N_{1}-N_{2})!} \frac{(M-n_{1}N_{1}-N_{2})!}{(M-n_{1}N_{1}-2N_{2})!}\cdots \frac{[M-n_{1}N_{1} - (n_{2}-1)N_{2}]!}{(M-n_{1}N_{1}-n_{2}N_{2})!}\\
&= \frac{1}{n_{2}!} \left[\frac{z-1}{M}\right]^{n_{2}(N_{2}-1)} \frac{(M-n_{1}N_{1})!}{(M-n_{1}N_{1}-n_{2}N_{2})!} \\
&= \frac{1}{n_{2}!} \left[\frac{z-1}{M}\right]^{n_{2}(N_{2}-1)} \frac{(M-n_{1}N_{1})!}{0!}
\end{aligned}
$$

\end{enumerate}

\end{itemize}

于是混合状态下的方法数为 

$$
\begin{aligned}
\Omega_{\text{mix}} &= W_{1}W_{2} \\
&= \frac{1}{n_{1}!}\left[\frac{z-1}{M}\right]^{n_{1}(N_{1}-1)}\frac{M!}{(M-n_{1}N_{1})!} \frac{1}{n_{2}!} \left[\frac{z-1}{M}\right]^{n_{2}(N_{2}-1)} \frac{(M-n_{1}N_{1})!}{0!}\\
&= \frac{M!}{n_{1}!n_{2}!}\left[\frac{z-1}{M}\right]^{n_{1}(N_{1}-1) + n_{2}(N_{2}-1)}
\end{aligned}
$$

混合熵被定义为 Boltzmann 熵变化量 $\begin{aligned}\Delta S = k_{B}\ln{\left(\frac{\Omega_{\text{mix}}}{\Omega_{\text{sep}}}\right)}\end{aligned}$, 而

$$
\frac{\Omega_{\text{mix}}}{\Omega_{\text{sep}}} = \frac{\frac{M!}{n_{1}!n_{2}!}\left[\frac{z-1}{M}\right]^{n_{1}(N_{1}-1) + n_{2}(N_{2}-1)}}{\frac{M_{1}!M_{2}!}{n_{1}!n_{2}!} \left[\frac{z-1}{M_{1}}\right]^{n_{1}(N_{1}-1)} \left[\frac{z-1}{M_{2}}\right]^{n_{2}(N_{2}-1)}} = \frac{M!}{M_{1}!M_{2}!}\frac{M_{1}^{n_{1}(N_{1}-1)}M_{2}^{n_{2}(N_{2}-1)}}{M^{n_{1}(N_{1}-1) + n_{2}(N_{2}-1)}}
$$

两边同取对数, 且使用 Stirling 近似 $\ln{X!}\approx X\ln{X} - X$, 第一项

$$
\begin{aligned}
\ln{\frac{M!}{M_{1}!M_{2}!}} &\approx M\ln{M} - M - M_{1}\ln{M_{1}} + M_{1} - M_{2}\ln{M_{2}} + M_{2} \\
&= M\ln{M} - M_{1}\ln{M_{1}} - M_{2}\ln{M_{2}}\\
&= (M_{1}+M_{2})\ln{M} - M_{1}\ln{M_{1}} - M_{2}\ln{M_{2}}\\
&= -M_{1}\ln{\left(\frac{M_{1}}{M}\right)} - M_{2}\ln{\left(\frac{M_{2}}{M}\right)}
\end{aligned}
$$

第二项

$$
\begin{aligned}
\ln{\left(\frac{M_{1}^{n_{1}(N_{1}-1)}}{M^{n_{1}(N_{1}-1)}}\frac{M_{2}^{n_{2}(N_{2}-1)}}{M^{n_{2}(N_{2}-1)}}\right)} &= n_{1}(N_{1}-1)\ln{\left(\frac{M_{1}}{M}\right)} + n_{2}(N_{2}-1)\ln{\left(\frac{M_{2}}{M}\right)}\\
&= (M_{1}-n_{1}) \ln{\left(\frac{M_{1}}{M}\right)} + (M_{2}-n_{2})\ln{\left(\frac{M_{2}}{M}\right)}
\end{aligned}
$$

约定体积分数为 $\phi_{i}$ 聚合物占据格点数 $N_{i}$ 与总格点数 $M$ 之比, 即 $\begin{aligned}\phi_{1} = \frac{M_{1}}{M} = \frac{n_{1}N_{1}}{M}\end{aligned}$, $\begin{aligned}\phi_{2} = \frac{M_{2}}{M} = \frac{n_{2}N_{2}}{M} = 1- \phi_{1}\end{aligned}$. 于是混合熵为

$$
\begin{aligned}
\frac{\Delta S_{\text{mix}}}{k_{B}} &= -M_{1}\ln{\phi_{1}} - M_{2}\ln{\phi_{2}} + (M_{1}-n_{1}) \ln{\phi_{1}} + (M_{2}-n_{2})\ln{\phi_{2}}\\
&= - n_{1}\ln{\phi_{1}} - n_{2}\ln{\phi_{2}}\\
&= M\left[ - \frac{\phi_{1}}{N_{1}}\ln{\phi_{1}} - \frac{\phi_{2}}{N_{2}}\ln{\phi_{2}} \right]
\end{aligned}
$$

单位格点的混合熵为

$$
\begin{aligned}
\frac{1}{M}\frac{\Delta S_{\text{mix}}}{k_{B}} = \boxed{\frac{\Delta\bar{S}_{\text{mix}}}{k_{B}} = -\frac{\phi_{1}}{N_{1}}\ln{\phi_{1}} - \frac{\phi_{2}}{N_{2}}\ln{\phi_{2}}}
\end{aligned}
$$

\begin{center}
\rule{0.8\textwidth}{0.4pt}
\end{center}

[Discussion] Flory-Huggins 理论推导

首先需要推导出 Flory-Huggins 自由能, 核心假设是溶剂和 polymer 的单体都会占据一个格点. 

因此设总格点数为 $M$, 各格点体积为 $a^{3}$, polymer 长度均为 $N$, 且一共有 $n_{p}$ 条聚合物链. 

每个溶剂分子占据 1 个格点, 且共有 $n_{s}$ 个溶剂分子. 因此存在关系式

$$
M =  Nn_{p} + n_{s} 
$$

polymer 的体积分数为 $\begin{aligned}\phi = \frac{Nn_{p}}{M}\end{aligned}$, 溶剂的体积分数为 $\begin{aligned}1 - \phi = \frac{n_{s}}{M}\end{aligned}$. 可以把溶剂分子视为聚合度为 1 的聚合物. 

自由能 $F \equiv U - TS$, 混合(视为等温过程)后产生变化 $\Delta F_{\text{mix}} = \Delta U_{\text{mix}} - T\Delta S_{\text{mix}}$. 也可通过同除以格点数 $M$ 和 能量量纲 $k_{B}T$ 化为形式 $\begin{aligned}\frac{\Delta\bar{F}_{\text{mix}}}{k_{B}T} = \frac{\Delta\bar{U}_{\text{mix}}}{k_{B}T} - \frac{\Delta\bar{S}_{\text{mix}}}{k_{B}}\end{aligned}$. 

设晶格配位数为 $z$. 那么将存在 $\varepsilon_{pp}$ (溶质-溶质), $\varepsilon_{ss}$ 溶剂-溶剂, $\varepsilon_{ps}$ 溶质-溶剂三种相互作用的能量. 

对于溶质-溶剂(polymer-solvent), 采取平均场假设, 即溶质和溶剂配位的概率为 $P = \phi(1-\phi)$. 则 polymer-solvent 的总配位数为 $\begin{aligned}N_{ps} = \frac{1}{2}zM\phi(1-\phi)\end{aligned}$, 其中 $\begin{aligned}\frac{1}{2}\end{aligned}$ 是为了消除重复计数. 

交换能被定义为 $\begin{aligned}\varepsilon = \varepsilon_{ps} - \frac{1}{2}(\varepsilon_{pp} + \varepsilon_{ss})\end{aligned}$. 那么由混合引起的内能变化则为

$$
\begin{aligned}
\Delta U_{\text{mix}} &= N_{ps}\varepsilon = \frac{z\varepsilon}{2}M\phi(1-\phi) \\
\Rightarrow \frac{\Delta\bar{U}_{\text{mix}}}{k_{B}T} &= \frac{\Delta U_{\text{mix}}}{Mk_{B}T} = \frac{z\varepsilon}{2k_{B}T} \phi (1-\phi) = \chi\phi(1-\phi)
\end{aligned}
$$

其中 $\begin{aligned}\frac{z\varepsilon}{2k_{B}T}\end{aligned}$ 被定义为 Flory-Huggins 参数. 

代入前面推导的混合熵公式, 即得到各格点自由能变化量(Flory-Huggins 自由能密度):

$$
\boxed{\frac{\Delta\bar{F}_{\text{mix}}}{k_{B}T} = \frac{\phi}{N}\ln{\phi} + (1-\phi)\ln{(1-\phi)} + \chi\phi(1-\phi)}
$$

在稀溶液极限 ($\phi\to 0$) 下, 有 $\begin{aligned}\ln{(1-x)} = - x - \frac{x^{2}}{2}\end{aligned}$, 则

$$
(1-\phi)\ln{(1-\phi)} = -\phi + \frac{\phi^{2}}{2} + \mathcal{O}(\phi^{3})
$$

自由能可写作

$$
\frac{\Delta\bar{F}_{\text{mix}}}{k_{B}T} = \overset{\text{ideal mix entropy}}{\frac{\phi}{N}\ln{\phi}} + \underset{A_{1}}{(\chi-1)}\phi + \underset{A_{2}}{\left(\frac{1}{2}-\chi\right)}\phi^{2} + \mathcal{O}(\phi^{3})
$$

这个形式来源于 Virial expansion: 

$$
\begin{aligned}
\frac{f}{k_{B}T} = \frac{\phi}{N}\ln{\phi} + A_{1}\phi + A_{2}\phi^{2} + A_{3}\phi^{3} + \cdots
\end{aligned}
$$

其中 $A_{1}$ 代表的是溶质和溶剂的平均能量贡献, $A_{2}$ 是溶质与溶质的两体相互作用. 

$$
A_{2}\sim \frac{v}{a^{3}}\Rightarrow v\propto a^{3}\left(\frac{1}{2}-\chi\right) \Rightarrow v = a^{3}(1-2\chi)
$$

这样就回到了熟悉的通过体积排斥参数 $v$ 来讨论溶剂性能的领域. 

\begin{itemize}
    \item 当 $T$ 较高, 使得 $\begin{aligned}\chi = \frac{z\varepsilon}{2k_{B}T}<\frac{1}{2}\end{aligned}$, 则 $v = a^{3}(1-2\chi)>0$, 链膨胀;
    \item 当 $\begin{aligned}T = T_{\theta} = \frac{z\varepsilon}{k_{B}}\end{aligned}$, 使得 $\begin{aligned}\chi = \frac{z\varepsilon}{2k_{B}T} = \frac{1}{2}\end{aligned}$, 则 $v = a^{3}(1-2\chi)=0$, 理想链;
    \item 当 $T$ 较低, 使得 $\begin{aligned}\chi = \frac{z\varepsilon}{2k_{B}T}>\frac{1}{2}\end{aligned}$, 则 $v = a^{3}(1-2\chi)<0$, 链塌缩;
\end{itemize}

根据 $T_{\theta}$ 的定义, 可将其代入至 $\chi$ 式中替换系数: 

$$
v = a^{3} \left(1 - 2 \frac{T_{\theta}}{2T}\right) = a^{3} \left(1 - \frac{T_{\theta}}{T}\right)\propto \boxed{1 - \frac{T_{\theta}}{T}}
$$

\begin{center}
\rule{0.8\textwidth}{0.4pt}
\end{center}}}

至于温度则是相当于在改变排斥体积参数 $v$, 因此可以视为温度在将溶剂在好溶剂在差溶剂之间调节. 体积排斥参数 $v$ 的温度 $T$ 依赖关系为 $\begin{aligned}v\propto 1 - \frac{T_{\theta}}{T}\end{aligned}$, 即低温 ($T < T_{\theta}$) 时体现为差溶剂, $T = T_{\theta}$ 时为 $\theta$ 溶剂, 高温 ($T > T_{\theta}$) 时为好溶剂.


\end{enumerate}

\subsection{液晶}

\begin{enumerate}
    \item \textbf{解释液晶的三种相: 丝状/向列相, 层状相和螺旋相}

液晶是取向有序, 位置长程无序的物质状态. 有序程度从高到低: 

\begin{figure}[htbp]
\centering
\includegraphics[width=0.4\textwidth]{image/lc-1.png}
\end{figure}

\begin{itemize}
    \item *Crystalline(结晶相): 条件最强的相, 要求分子取向完全一致, 位置如同晶格一样严格有序;
    \item Smectic(层状相): 分子长轴取向一致, 相比于 Nematic 具有层状结构, 存在一维位置有序. 层和层之间可以滑动;
\begin{figure}[htbp]
\centering
\includegraphics[width=0.4\textwidth]{image/lc-2.png}
\end{figure}
    \item Nematic(向列相): 分子长轴取向一致, 无层状结构, 位置(液晶分子中心)长程无序. 被比喻为 "盒子里的牙签", 总体方向一样但是可以自由向任意位置移动;
\begin{figure}[htbp]
\centering
\includegraphics[width=0.4\textwidth]{image/lc-3.png}
\end{figure}
    \item Cholesteric(螺旋相): Nematic 的 director $\vec{n}$ 沿着空间中某轴呈现螺旋;
\begin{figure}[htbp]
\centering
\includegraphics[width=0.4\textwidth]{image/lc-4.png}
\end{figure}
    \item Isotropic(各向同性相): 分子取向和位置完全无序, 是液晶在高温下表现出的液体性质, 物理性质呈现各向同性. $S = 0$. 
\end{itemize}

    \item \textbf{解释 splay, twist 和 bend 三种自由能}

director 在空间中产生变化时, 即有弹性自由能, 可被分为三个部分: splay(扩张), twist(扭转) 和 bend(弯曲). 

$$
f = \underset{\text{Splay}}{\frac{1}{2}K_{1}(\nabla\cdot\vec{n})^{2}} + \underset{\text{Twist}}{\frac{1}{2}K_{2}(\vec{n}\cdot\nabla\times\vec{n})^{2}} + \underset{\text{Bend}}{\frac{1}{2}K_{3}|\vec{n}\times\nabla\times\vec{n}|^{2}}
$$

若材料近似各向同性($K_{1}\approx K_{2}\approx K_{3} = K$), 则

$$
f = \frac{1}{2}K(\nabla \vec{n})^{2}
$$

    \item \textbf{解释液晶的序参量定义}
    
{\color{gray}{[Discussion] 为何液晶需要新的序参量定义. 

假设液晶分子长轴取向为单位向量 $\vec{u}$. 于是平均取向可以被定义为 $\vec{n} = \left\langle \vec{u} \right\rangle$. $\vec{u}$ 和 $\vec{n}$ 之间的夹角 $\theta$ 可以被视为描述液晶分子的参数. 接下来需要定义描述液晶分子的序参量, 即其大小和液晶整体的有序性相关. 

在 Ising Model 中, 序参量被定义为 $\begin{aligned}L = \frac{1}{N}\sum_{i}\sigma_{i} = \frac{N_{+}-N_{-}}{N}\end{aligned}$, 即自旋的平均值. 

\begin{itemize}
    \item 如果借用 Ising Model 的定义, 我们有 $S_{1} = \langle\theta\rangle$. 然而, 液晶分子通常自身具有对称性(即长轴取向 $\vec{n}$ 和 $-\vec{n}$ 等价), 偏离平均取向 $\vec{n}$ 的分子 $\theta$ 和 $\theta+\pi$ 实际上是等价的, $S_{1}$ 却会错误的计算出 $\langle\theta\rangle = \theta+\pi/2$ 的结果. 
    \item 序参量 $S_{2} = \left\langle\cos{\theta}\right\rangle$ 同样没好到哪里去, 虽然 $\cos{\theta}$ 具有 $\cos{(2\pi-\theta)} = \cos{\theta}$ 的对称性, 但是仍然具有 $\cos{(\theta+\pi)} = -\cos{\theta}$ 的问题: 明明取向相同, 却为 $S_{2}$ 添加了 $0$ 的成分. 
    \item 最简单的消去符号的方法是对其做平方, 即拟定序参量为 $S_{3} = \langle\cos^{2}{\theta}\rangle$. 我们注意到: 
    \begin{itemize}
        \item 完全有序即 $\theta = 0$ 时, $S_{3} = \cos^{2}{0} = 1$; 
        \item 完全无序时, $\theta$ 均匀分布在整个球壳上, 即 $\begin{aligned}
            f(\theta,\phi) = \frac{1}{4\pi}
        \end{aligned}$. 

    计算拟定序参量

    $$
    \begin{aligned}
    \langle\cos^{2}(\theta)\rangle &= \int_{0}^{2\pi}\mathrm{d}\phi \int_{0}^{\pi}\cos^{2}{\theta}\cdot \frac{1}{4\pi}\sin{\theta}\mathrm{d}\theta
    = -\frac{1}{2}\int_{0}^{\pi}\cos^{2}{\theta}\mathrm{d}(\cos{\theta}) \\
    &= -\frac{1}{2}\int_{1}^{-1}u^{2}\mathrm{d}u = \frac{1}{2}\frac{1}{3}u^{3}\bigg|_{-1}^{1} = \frac{1}{3}
    \end{aligned}
    $$

    我们希望在完全无序时, 序参量为 $0$. 因此对该单调函数进行线性变换即可得到最后结果. 
    \end{itemize}
\end{itemize}

从更深层次的角度而言, 任何方向分布函数都可以展开为 Legendre 多项式级数:

$$
f(\theta) = \sum_{l=0}^{\infty} \frac{2l+1}{2}\langle P_{l}(\cos{\theta})\rangle P_{l}(\cos{\theta})
$$

\begin{itemize}
    \item $l = 1$ 对应 $\langle\cos{\theta}\rangle$, 对应极性有序;
    \item $l = 2$ 对应 $\begin{aligned}
        \left\langle\frac{3\cos^{2}{\theta}-1}{2}\right\rangle
    \end{aligned}$, 对应取向有序. 
\end{itemize}}}

\begin{center}
\rule{0.8\textwidth}{0.4pt}
\end{center}

$$
S = \left\langle \frac{3\cos^{2}{\theta} - 1}{2} \right\rangle
$$

\begin{enumerate}
    \item 若 $S = 1$, 即所有 $\theta = 0$, 液晶分子取向完全平行一致(向列相及以上);
    \item 若 $S = 0$, 即 $\begin{aligned}
        \langle\cos^{2}{\theta}\rangle = \frac{1}{3}
    \end{aligned}$, 这是液晶分子完全无序计算的结果.
\end{enumerate}

    \item \textbf{解释概念 "液晶的 $\begin{aligned}
        \pm\frac{1}{2}
    \end{aligned}$ 缺陷"}

[Discussion] winding number 的定义. 

可以借助电荷的电场线来理解 winding number. 对于点电荷 $q$, 在二维平面内研究其围线上电场方向的旋转. 如果以顺时针方向沿着围线旋转一周, 会发现电场方向也是以顺时针方向旋转一周 ($2\pi$). 因此其拓扑荷为 $m = 2\pi / (2\pi) = 1$, 这一点对于正负电荷都一样. 

\begin{figure}[htbp]
\centering
\includegraphics[width=0.4\textwidth]{image/winding_number.png}
\end{figure}

拓扑荷被定义为

$$
\boxed{m = \frac{1}{2\pi}\oint \mathrm{d}\theta}
$$

而在液晶中, 特殊的是 $\vec{u}$ 和 $-\vec{u}$ 是等价的. 因此若沿围线一周后, 分子取向 $\vec{u}$ 旋转了 $\pm\pi$, 实际上就已经完成了取向旋转 "一周". 此时围线中将会出现 $\pm \frac{1}{2}$ 的缺陷核. 两者分别被戏称为 "彗星"[$\begin{aligned}
    \theta(\phi) = \frac{\phi}{2}
\end{aligned}$] 和 "三叶"[$\begin{aligned}
    \theta(\phi) = -\frac{\phi}{2}
\end{aligned}$]. 

\begin{figure}[htbp]
\centering
\includegraphics[width=0.5\textwidth]{image/defects.png}
\end{figure}

\end{enumerate}

\end{document}